% Original: PDF strana 10, donji slajd; prvi deo.
\slajdid{03-s020}
\begin{frame}[t,label=03-s020]{Dozvola kodovanja signalom EI}
\setlength{\parskip}{0pt}
\begin{columns}[c,onlytextwidth]
\begin{column}{.31\textwidth}
% element: 03-s020:e01
Ulaz \textcolor{DEisticanje}{\textbf{EI}} (\emph{enable input})
dozvoljava kodovanje:
\begin{itemize}
\item za \(EI=1\), kodovanje je dozvoljeno;
\item za \(EI=0\), izlazi \(A_1,A_0\) postavljaju se na \textbf{logičku nulu}.
\end{itemize}
\par\vspace{1.5mm}
Možemo usloviti i \(GS\). Tada \(GS=0\) znači da nema aktivnih ulaza
ili da je komponenti \textbf{zabranjeno kodovanje} zbog višeg prioriteta.
\end{column}
\begin{column}{.66\textwidth}\centering
% element: 03-s020:e02
\begingroup
\def\DEenable{1}\def\DEjedinica{.75cm}\def\DEvisinaPrioriteta{.70cm}\def\DEgejt{.55}
% DEenable: dodatni EI/EO; DEjedinica i DEgejt određuju geometriju, font ostaje 9 pt.
\providecommand{\DEvisinaPrioriteta}{\DEjedinica}%
\begin{tikzpicture}[circuit logic US,x=\DEjedinica,y=\DEvisinaPrioriteta,font=\fontsize{9}{11}\selectfont,
 inv/.style={not gate US,draw,minimum width=8mm,minimum height=7mm,scale=\DEgejt},
 kapija/.style={and gate US,draw,minimum width=10mm,minimum height=10mm,scale=\DEgejt},
 zbir/.style={or gate US,draw,minimum width=10mm,minimum height=12mm,scale=\DEgejt,rotate=-90}]
\node[kapija,logic gate inputs=nn,anchor=input 2] (p2) at (2.7,3.2) {};
\node[kapija,logic gate inputs=nnn,anchor=input 3] (p1) at (2.7,2) {};
\node[kapija,logic gate inputs=nnnn,anchor=input 4] (p0) at (2.7,.8) {};
\draw[DEvod] (-.4,4.4) node[left] {I3} -- (3.4,4.4);
\coordinate (p3) at (3.4,4.4);
\foreach \n/\y/\pin in {2/3.2/2,1/2/3,0/.8/4}
 \draw[DEvod] (-.4,\y) node[left] {I\n} --(p\n.input \pin);
\foreach \n/\x/\y/\src in {3/.9/3.9/4.4,2/.45/2.7/3.2,1/0/1.5/2}{
 \node[inv,rotate=-90] (i\n) at (\x,\y) {};
 \draw[DEvod] (\x,\src)--(i\n.input);
 \node[junction] at (\x,\src) {};
}
\coordinate (b3) at ($(i3.output)+(.18,0)$);
\coordinate (b2) at ($(i2.output)+(.18,0)$);
\draw[DEvod] (i3.output)--(b3);
\draw[DEvod] (i2.output)--(b2);
\foreach \n in {2,1,0}
 \draw[DEvod] (b3)|-(p\n.input 1);
\foreach \n in {1,0}
 \draw[DEvod] (b2)|-(p\n.input 2);
\draw[DEvod] (i1.output)|-(p0.input 3);
\node[junction] at (b3 |- p2.input 1) {};
\node[junction] at (b3 |- p1.input 1) {};
\node[junction] at (b2 |- p1.input 2) {};
\foreach \n in {2,1,0} \coordinate (p\n) at (p\n.output);
\ifnum\DEenable=1
 \foreach \n in {3,2,1,0}{
  \node[kapija,logic gate inputs=nn,anchor=input 2] (e\n) at (4.5,0 |- p\n) {};
  \draw[DEvod] (p\n)--(e\n.input 2);
  \draw[DEvod] (e\n.input 1)--(4.03,0 |- e\n.input 1);
  \node[junction] at (4.03,0 |- e\n.input 1) {};
  \coordinate (q\n) at (e\n.output);
 }
 \draw[DEvod] (-.4,5.05) node[left] {EI} --(4.03,5.05)--(4.03,0 |- e0.input 1);
 \node[junction] at (4.03,5.05) {};
 \def\DXone{6.1}\def\DXzero{7.5}\def\DXgs{9.0}\def\DXend{9.45}
\else
 \foreach \n in {3,2,1,0} \coordinate (q\n) at (p\n);
 \def\DXone{4.5}\def\DXzero{5.7}\def\DXgs{7.1}\def\DXend{7.7}
 \node at (1.45,5.15) {Prioritetna mreža};
 \node at (5.9,5.15) {Koderska mreža};
\fi
\foreach \n in {3,2,1,0} \draw[DEvod] (q\n)--(\DXend,0 |- q\n);
\node[zbir,logic gate inputs=nn] (a1) at (\DXone,-.25) {};
\node[zbir,logic gate inputs=nn] (a0) at (\DXzero,-.25) {};
\node[zbir,logic gate inputs=nnnn] (gs) at (\DXgs,-.25) {};
\foreach \pin/\n in {2/3,1/2}
 {\draw[DEvod] (a1.input \pin)--(a1.input \pin |- q\n);
  \node[junction] at (a1.input \pin |- q\n) {};}
\foreach \pin/\n in {2/3,1/1}
 {\draw[DEvod] (a0.input \pin)--(a0.input \pin |- q\n);
  \node[junction] at (a0.input \pin |- q\n) {};}
\foreach \pin/\n in {4/3,3/2,2/1,1/0}
 {\draw[DEvod] (gs.input \pin)--(gs.input \pin |- q\n);
  \node[junction] at (gs.input \pin |- q\n) {};}
\foreach \name/\lab in {a1/A1,a0/A0,gs/GS}
 \draw[DEvod] (\name.output)--(\name.output |- 0,-1.3) node[below] {\lab};
\ifnum\DEenable=1
 \node[junction] at (gs.output |- 0,-1.1) {};
 \node[inv,rotate=90] (ngs) at (10,-.35) {};
 \draw[DEvod] (gs.output)--(gs.output |- 0,-1.1)--(10,-1.1)--(ngs.input);
 \node[kapija,rotate=-90,logic gate inputs=nn] (eo) at (11.15,-.35) {};
 \draw[DEvod] (ngs.output)--(ngs.output |- 0,.35)--(eo.input 2 |- 0,.35)--(eo.input 2);
 \draw[DEvod] (4.03,5.05)--(4.03,5.5)--(11.6,5.5)--(11.6,.35)--(eo.input 1 |- 0,.35)--(eo.input 1);
 \draw[DEvod] (eo.output)--(eo.output |- 0,-1.3) node[below] {EO};
\fi
\end{tikzpicture}

\endgroup
\end{column}
\end{columns}
\end{frame}
